\documentclass[11pt, a4paper]{article}

% === Encoding e Font ===
\usepackage[utf8]{inputenc}
\usepackage[T1]{fontenc}
\usepackage{lmodern}

% === Lingua ===
\usepackage[italian]{babel}

% === Layout Pagina ===
\usepackage{geometry}
\geometry{
  top=2.5cm,
  bottom=2.5cm,
  left=2.5cm,
  right=2.5cm,
}

% === Microtipografia ===
\usepackage{microtype}

% === Float ===
\usepackage{float}

% === Tabelle ===
\usepackage{array}
\usepackage{booktabs}
\usepackage{longtable}
\usepackage{multirow}
\usepackage{tabularx}

% === Colori ===
\usepackage[dvipsnames,svgnames,x11names]{xcolor}

% === Liste ===
\usepackage{enumitem}

% === Hyperlink ===
\usepackage{hyperref}
\hypersetup{
  colorlinks=true,
  linkcolor=blue!70!black,
  urlcolor=blue!70!black,
  citecolor=green!50!black,
  pdfauthor={WordToLaTeX},
  pdftitle={Documento di Test},
}

% === Utilità ===
\usepackage{parskip}       % Spaziatura tra paragrafi
\usepackage{setspace}       % Interlinea
\usepackage{fancyhdr}       % Header/Footer
\usepackage{lastpage}       % Numero ultima pagina
\usepackage{caption}        % Caption per figure/tabelle
\usepackage{textcomp}       % Simboli aggiuntivi

% === Header e Footer ===
\pagestyle{fancy}
\fancyhf{}
\fancyhead[L]{\small Documento di Test}
\fancyhead[R]{\small \thepage\ / \pageref{LastPage}}
\renewcommand{\headrulewidth}{0.4pt}
\renewcommand{\footrulewidth}{0pt}

% === Interlinea ===
\onehalfspacing

\begin{document}

\title{Documento di Test}
\author{WordToLaTeX}
\date{}
\maketitle
\thispagestyle{fancy}

Questo è un documento di prova per testare la conversione da \textbf{Word a PDF} tramite \textit{LaTeX}. Il convertitore gestisce \textbf{grassetto}, \textit{corsivo}, \underline{sottolineato} e anche \textcolor[RGB]{255,0,0}{testo colorato}.

\section{Sezione 1: Formattazione del Testo}

Il testo può essere formattato in diversi modi:

\textit{\textbf{Grassetto e Corsivo insieme}}

Testo con dimensione grande

Apice: E=mc\textsuperscript{2} e Pedice: H\textsubscript{2}O

\section{Sezione 2: Liste}

\begin{itemize}
  \item Elemento puntato 1
  \item Elemento puntato 2
  \item Elemento puntato 3
\end{itemize}

\begin{enumerate}
  \item Primo elemento numerato
  \item Secondo elemento numerato
  \item Terzo elemento numerato
\end{enumerate}

\section{Sezione 3: Tabelle}

\begin{table}[H]
  \centering
  \begin{tabularx}{\textwidth}{|X|X|X|}
    \hline
    \textbf{Nome} & \textbf{Cognome} & \textbf{Età} \\
    \hline
    Mario & Rossi & 35 \\
    \hline
    Lucia & Bianchi & 28 \\
    \hline
    Andrea & Verdi & 42 \\
    \hline
  \end{tabularx}
\end{table}

\section{Sezione 4: Struttura Gerarchica}

\subsection{Sottosezione 4.1}

Contenuto della sottosezione 4.1.

\subsection{Sottosezione 4.2}

Contenuto della sottosezione 4.2.

\subsubsection{Sotto-sottosezione 4.2.1}

Contenuto più profondo nella gerarchia.

\section{Sezione 5: Allineamento}

Testo allineato a sinistra.

\begin{center}
Testo centrato.
\end{center}

\begin{flushright}
Testo allineato a destra.
\end{flushright}

Testo giustificato: Lorem ipsum dolor sit amet, consectetur adipiscing elit. Sed do eiusmod tempor incididunt ut labore et dolore magna aliqua. Ut enim ad minim veniam, quis nostrud exercitation ullamco laboris nisi ut aliquip ex ea commodo consequat.


\end{document}